\documentclass[11pt,a4paper]{article}

\usepackage[T1]{fontenc}
\usepackage[utf8]{inputenc}
\usepackage[italian]{babel}
\usepackage{lmodern}
\usepackage{microtype}
\usepackage[a4paper,margin=2.15cm,headheight=15pt]{geometry}
\usepackage{booktabs,longtable,tabularx,array}
\usepackage{ragged2e}
\usepackage{enumitem}
\usepackage{xcolor}
\usepackage[most]{tcolorbox}
\usepackage{fancyhdr}
\usepackage{hyperref}
\usepackage{xurl}
\usepackage{listings}
\usepackage{titlesec}
\usepackage{amsmath,amssymb}

\definecolor{ddtadark}{RGB}{28,38,52}
\definecolor{ddtagray}{RGB}{244,246,248}
\definecolor{ddtablue}{RGB}{42,88,135}
\definecolor{ddtagreen}{RGB}{48,111,78}
\definecolor{ddtaorange}{RGB}{148,91,25}
\definecolor{ddtared}{RGB}{140,48,48}
\definecolor{ddtapurple}{RGB}{92,63,122}

\hypersetup{
  colorlinks=true,
  linkcolor=ddtablue,
  urlcolor=ddtablue,
  citecolor=ddtablue,
  pdftitle={Base Analysis Operational Guide},
  pdfauthor={Documentation-Driven Threat Analysis research workspace}
}

\pagestyle{fancy}
\fancyhf{}
\lhead{DDTA -- Base Analysis Operational Guide}
\rhead{Revision 2 - BA0--BA6 integrated method}
\cfoot{\thepage}

\titleformat{\section}{\Large\bfseries\color{ddtadark}}{\thesection}{0.7em}{}
\titleformat{\subsection}{\large\bfseries\color{ddtadark}}{\thesubsection}{0.7em}{}
\titleformat{\subsubsection}{\normalsize\bfseries\color{ddtadark}}{\thesubsubsection}{0.7em}{}

\setlist[itemize]{leftmargin=1.55em,itemsep=0.22em,topsep=0.35em}
\setlist[enumerate]{leftmargin=1.75em,itemsep=0.22em,topsep=0.35em}
\setlength{\parskip}{0.35em}
\setlength{\parindent}{0pt}
\setlength{\emergencystretch}{3em}

\lstset{
  basicstyle=\ttfamily\small,
  columns=fullflexible,
  keepspaces=true,
  frame=single,
  framerule=0.3pt,
  rulecolor=\color{ddtablue},
  backgroundcolor=\color{ddtagray},
  xleftmargin=0.6em,
  xrightmargin=0.6em,
  aboveskip=0.5em,
  belowskip=0.5em,
  breaklines=true,
  showstringspaces=false
}

\newtcolorbox{closedbox}[1][REGOLA STABILE]{enhanced,breakable,colback=ddtagray,colframe=ddtagreen,title={#1},fonttitle=\bfseries}
\newtcolorbox{workingbox}[1][REGOLA BASELINE-SCOPED]{enhanced,breakable,colback=white,colframe=ddtaorange,title={#1},fonttitle=\bfseries}
\newtcolorbox{openbox}[1][NOT SPECIFIED / DIAGNOSTIC]{enhanced,breakable,colback=white,colframe=ddtapurple,title={#1},fonttitle=\bfseries}
\newtcolorbox{goodbox}[1][PREFERIBILE]{enhanced,breakable,colback=white,colframe=ddtagreen,title={#1},fonttitle=\bfseries}
\newtcolorbox{badbox}[1][DA EVITARE]{enhanced,breakable,colback=white,colframe=ddtared,title={#1},fonttitle=\bfseries}
\newtcolorbox{examplebox}[1][ESEMPIO FACIAL ACCESS]{enhanced,breakable,colback=ddtagray,colframe=ddtablue,title={#1},fonttitle=\bfseries}
\newtcolorbox{conceptbox}[1][CONCETTO]{enhanced,breakable,colback=white,colframe=ddtablue,title={#1},fonttitle=\bfseries}

\newcolumntype{Y}{>{\RaggedRight\arraybackslash}X}
\newcolumntype{L}[1]{>{\RaggedRight\arraybackslash}p{#1}}

\newcommand{\code}[1]{\nolinkurl{#1}}
\newcommand{\BA}{\textit{Base Analysis}}
\newcommand{\BAR}{\textit{BAReferent}}
\newcommand{\BAP}{\textit{BAProposition}}

\begin{document}

\begin{titlepage}
\centering
\vspace*{1.5cm}
{\Large Documentation-Driven Threat Analysis\par}
\vspace{0.8cm}
{\Huge\bfseries Base Analysis Operational Guide\par}
\vspace{0.25cm}
{\LARGE Revision 2 - BA0--BA6 integrated method\par}
\vspace{0.9cm}
{\large Procedura source-first dalla baseline governata alla accepted Base Analysis e alle projection downstream\par}
\vfill
\begin{minipage}{0.88\textwidth}
\small
\textbf{Stato.} Consolidamento human-readable successivo alla chiusura BA6 sul corpus governato Facial Access.

\medskip
\textbf{Project authority.} \code{FACIAL-ACCESS-GOV-R2}.

\medskip
\textbf{Accepted BA baseline.} \code{FACIAL-ACCESS-BA-R24-R1}.

\medskip
\textbf{Pinned source revision.} \code{8af2257a1df94fa5a83d4853ed0a1eb4d020c429}.

\medskip
Questa guida non sostituisce i contratti normativi BA0--BA6 né i documenti governati del progetto. Li organizza in una procedura leggibile e operativa.
\end{minipage}
\vfill
{\large 29 agosto 2026\par}
\end{titlepage}

\pagenumbering{roman}
\tableofcontents
\newpage
\pagenumbering{arabic}

\section{Che cos'è la Base Analysis}

Per una baseline documentale governata, la Base Analysis è la rappresentazione analitica condivisa, methodology-neutral, del significato di progetto necessario a DDTA.

\begin{closedbox}[Authority boundary]
\begin{lstlisting}
governed project documentation
    = project authority

accepted Base Analysis
    = shared analytical baseline

projection / threat method / tests
    = downstream consumers
\end{lstlisting}
\end{closedbox}

La BA può normalizzare e rendere riusabile il significato, ma non può creare silenziosamente project commitment.

\section{Il contratto BA0--BA6}

La metodologia corrente separa sette responsabilità.

\begin{longtable}{L{0.13\textwidth}L{0.75\textwidth}}
\toprule
\textbf{Layer} & \textbf{Responsabilità} \\
\midrule\endhead
BA0 & authority boundary, minimalità, unresolved state, feedback direction. \\
BA1 & due sole famiglie first-class: \code{BAReferent} e \code{BAProposition}. \\
BA2 & proposition semantics, operatori, ruoli, cardinalità e strutture locali. \\
BA3 & provenance, derivazione, review/freshness, continuity e revalidation. \\
BA4 & projection, traceability, interpretation e coverage downstream. \\
BA5 & canonical semantic registries e controlled authoring. \\
BA6 & accettazione integrata di una concreta BA baseline dopo coverage e regression. \\
\bottomrule
\end{longtable}

BA6 non rende BA0--BA5 ``più veri'': verifica che una concreta materializzazione li soddisfi insieme.

\section{BAReferent e BAProposition}

\subsection{BAReferent}

Un \BAR{} è un significato di progetto methodology-neutral con identità indipendente quando deve essere:

\begin{itemize}
\item riusato da più proposition;
\item qualificato o vincolato;
\item correlato con altri significati;
\item confrontato tra baseline;
\item target di più asserzioni;
\item tracciato stabilmente attraverso change/revalidation.
\end{itemize}

Può denotare participant, capability, information, behavior, state, service, contract, boundary o altro significato. Queste categorie non sono automaticamente nuove famiglie BA.

\subsection{BAProposition}

Una \BAP{} è una asserzione analitica su uno o più BAReferent.

\begin{closedbox}[No proposition proxy]
La proposition non è il comportamento o l'oggetto di progetto che descrive. Se quel significato deve essere riusato indipendentemente, riceve BAReferent identity.
\end{closedbox}

Questa regola è ciò che giustifica \code{RecognitionCaptureDelivery} come BAReferent separato dalla proposition \code{transfer} che lo descrive.

\section{Il criterio di dettaglio}

La BA non replica l'intera documentazione e non ricostruisce automaticamente una architettura completa.

\begin{lstlisting}
more detail
    only if needed to:
      preserve a governed distinction
      answer a material downstream question
      or expose a missing governed fact
\end{lstlisting}

Stopping test:

\begin{conceptbox}
Abbiamo abbastanza semantica per formulare la domanda e distinguere evidenza supportata, contraddetta o \code{NOT SPECIFIED}?
\end{conceptbox}

Se sì, stop. Se no, si identifica la distinzione mancante o si produce una diagnosis. Non si inventa l'implementazione.

\section{Procedura operativa}

\begin{enumerate}
\item Pin della source authority.
\item Lettura source-first.
\item Identificazione del minimum shared meaning.
\item Creazione dei BAReferent necessari.
\item Formulazione delle BAProposition con BA2.
\item Collegamento BA3 alla source evidence.
\item Registrazione di ambiguity / NOT SPECIFIED.
\item Minimum-BA acceptance preliminare.
\item Regression completa verso le fonti.
\item Pressure review della metodologia solo su counterexample concreti.
\item BA6 integrated completion.
\item Projection/downstream consumption secondo BA4.
\end{enumerate}

\section{BA2 R3 - operatori}

La revisione R3 mantiene quattordici operatori.

\begin{longtable}{L{0.22\textwidth}L{0.66\textwidth}}
\toprule
\textbf{Operator} & \textbf{Significato} \\
\midrule\endhead
\code{transfer} & conveyance di content da source a destination; optional behavior identity quando il trasferimento deve essere riusato/qualificato. \\
\code{produce} & actor/capability rende disponibile un risultato governato. \\
\code{create} & stabilisce un nuovo item/event project-semantic. \\
\code{observe} & legge/ispeziona significato esistente senza asserire creazione o state change. \\
\code{transition} & cambia stato/lifecycle di un subject. \\
\code{correlate} & preserva same-request/evaluation/identity/context binding. \\
\code{reference} & riferimento esplicito senza implicare dependency/correlation/ownership. \\
\code{dependOn} & prerequisite o dipendenza. \\
\code{consumeService} & consumo effettivo di servizio senza ownership transfer. \\
\code{realize} & significato concreto realizza/materializza significato astratto. \\
\code{assignResponsibility} & placement o negazione di responsibility/authority kind. \\
\code{constrain} & restriction riusabile/queryable. \\
\code{classify} & semantic kind method-neutral. \\
\code{decisionRule} & mapping operativo da condizioni/input a result. \\
\bottomrule
\end{longtable}

\section{Ruoli e cardinalità principali}

\begin{longtable}{L{0.19\textwidth}L{0.33\textwidth}L{0.34\textwidth}}
\toprule
\textbf{Operator} & \textbf{Required} & \textbf{Optional} \\
\midrule\endhead
transfer & source 1; destination 1..*; content 1..* & behavior 0..1 \\
produce & actor 1; result 1..* & input 0..* \\
create & actor 1; created 1..* & - \\
observe & actor 1; observed 1..* & result 0..* \\
transition & subject 1; toState 1 & actor 0..1; fromState 0..1 \\
correlate & correlatedItem 1..*; correlationContext 1 & - \\
reference & referencer 1; referenced 1..* & - \\
dependOn & dependent 1; prerequisite 1..* & - \\
consumeService & consumer 1..*; service 1 & provider 0..1 \\
realize & abstract 1; realization 1..* & - \\
assignResponsibility & responsibleParty 1..*; responsibilityScope 1; responsibilityKind 1 & - \\
constrain & constraintTarget 1; constraintValue 1..* & - \\
classify & classifiedReferent 1; semanticKind 1..* & - \\
decisionRule & actor 1; input 1..*; result 1 & operator-local rule \\
\bottomrule
\end{longtable}

\section{transfer.behavior}

R3 introduce una sola estensione al transfer:

\begin{lstlisting}
transfer
  behavior    -> BAReferent [0..1]
  source      -> BAReferent [1]
  destination -> BAReferent [1..*]
  content     -> BAReferent [1..*]
\end{lstlisting}

Usare \code{behavior} solo quando il comportamento di trasferimento deve essere target di altre proposition.

\begin{examplebox}
\begin{lstlisting}
transfer
  behavior    -> RecognitionCaptureDelivery
  source      -> CameraSubsystem
  destination -> RecognitionProcessor
  content     -> RecognitionCapture
\end{lstlisting}

\code{RecognitionCaptureDelivery} permette di mantenere sullo stesso segmento:
\begin{itemize}
\item ConnectivityService dependency;
\item WIRED\_ETHERNET;
\item Confidentiality;
\item Integrity;
\item AuthorizedProvenance.
\end{itemize}
\end{examplebox}

Non si generalizza \code{behavior} a tutti gli operatori finché un secondo counterexample non lo richiede.

\section{decisionRule}

Forma minima:

\begin{lstlisting}
decisionRule
  actor  -> BAReferent
  input  -> BAReferent [1..*]
  result -> BAReferent

  rule
    IF    decisionCondition
    THEN  resultAssignment [1..*]
    ELSE  resultAssignment [0..*]
\end{lstlisting}

Condition forms ammesse:

\begin{lstlisting}
comparison
satisfies
allOf
anyOf
not
\end{lstlisting}

\subsection{comparison}

\begin{lstlisting}
comparison
  referent      -> BAReferent
  property      -> controlled semantic key
  comparisonKey -> equals | notEquals
  value         -> typed local value | BAReferent
\end{lstlisting}

\code{property} resta obbligatoria.

\subsection{satisfies}

\begin{lstlisting}
satisfies
  subject   -> BAReferent
  condition -> BAReferent
\end{lstlisting}

È una forma locale a \code{decisionRule}. Serve quando la source governa una condizione ma non ne governa il modello interno property/value.

\begin{examplebox}[FR-1.1]
\begin{lstlisting}
allOf
  comparison
    IdentityDeterminationOutcome.outcomeKind == SUCCESS

  satisfies
    subject   -> AccessAuthorizationState
    condition -> RequiredAccessAuthorizationCondition
\end{lstlisting}

Questo preserva ``authorization condition satisfied'' senza inventare:

\begin{lstlisting}
AccessAuthorizationState.authorized = TRUE
\end{lstlisting}
\end{examplebox}

\section{constrain e property vocabulary}

Quando la documentazione governa un dominio di valori ma non la regola di selezione:

\begin{lstlisting}
constrain
  constraintTarget -> IdentityDeterminationOutcome
  constraintValue
    property   -> outcomeKind
    vocabulary -> [SUCCESS, NEGATIVE, INCONCLUSIVE]
\end{lstlisting}

Non usare \code{constrain} come escape hatch per prose non strutturate.

\section{correlate}

\code{correlate} preserva binding di contesto.

Esempi:

\begin{lstlisting}
RecognitionCapture
    <-> IdentityDeterminationRequest

IdentityDeterminationOutcome
AccessAuthorizationState
    <-> same GovernedIdentity

IdentityDeterminationOutcome
    <-> AccessAttempt
\end{lstlisting}

Non usare \code{reference} come sostituto quando il requisito richiede identità di contesto.

\section{Acquisition: quando non scegliere un operatore}

Il verbo documentale \code{acquire} non forza un nuovo operatore.

\begin{lstlisting}
create
observe
produce
\end{lstlisting}

sono interpretazioni materialmente diverse.

La BA accettata mantiene:

\begin{lstlisting}
RecognitionCaptureAcquisition [BAReferent]
diagnosis:
    MULTIPLE MATERIAL BA CANDIDATES
\end{lstlisting}

e non materializza una proposition falsa solo per avere un arco nel grafo.

\section{BA3 - provenance}

Ogni BA element meaning-bearing deve risolvere almeno a:

\begin{lstlisting}
governedBaselineKey
originState
sourceLink
\end{lstlisting}

Origin states:

\begin{lstlisting}
GROUNDED
DERIVED
DIAGNOSTIC_UNRESOLVED
\end{lstlisting}

\subsection{sourceLink non è semantica BA2}

Tre responsabilità restano distinte:

\begin{lstlisting}
sourceLink
derivationBasisBinding
revalidationContext
\end{lstlisting}

Non usare provenance per simulare una relazione che BA2 non sa rappresentare.

\section{Derivazione}

Un elemento \code{DERIVED} richiede:

\begin{itemize}
\item role-bound derivation basis;
\item immutable derivation rule revision;
\item applicability contract;
\item output semantic contract;
\item lineage fino a governed source.
\end{itemize}

Nel baseline Facial Access corrente le proposition accettate sono normalizzazioni GROUNDED; non è necessario inventare derived semantics per completare la BA.

\section{Diagnostics}

La BA può accettare una diagnosis come diagnosis.

Vocabolario utile:

\begin{lstlisting}
MULTIPLE MATERIAL BA CANDIDATES
BA REQUIRES UNSUPPORTED FACT
BA LOSES GOVERNED MATERIAL DISTINCTION
BA EXPOSES SOURCE CONFLICT
NOT SPECIFIED
\end{lstlisting}

\begin{closedbox}
Una diagnosis non è una project truth.
\end{closedbox}

\section{BA5 - authoring canonico}

I semantic keys operativi devono essere canonici.

Domini distinti:

\begin{itemize}
\item project referent names;
\item BA2 operator keys;
\item operator-scoped role keys;
\item semantic kind keys;
\item controlled value domains.
\end{itemize}

Un synonym o una traduzione non stabiliscono equivalenza normativa automaticamente.

\section{BA4 - projection boundary}

Una projection è downstream.

Catena:

\begin{lstlisting}
governed documentation
    -> accepted BA
        -> projection
            -> method-owned interpretation
\end{lstlisting}

Ogni meaning-bearing projection item deve tracciare alla BA.

Coverage modes:

\begin{lstlisting}
EXHAUSTIVE_FOR_DECLARED_SCOPE
SELECTIVE
\end{lstlisting}

Una selective projection può omettere elementi BA senza dichiararli falsi o assenti.

\section{BA6 - integrated completion}

BA6 accetta una concreta baseline analitica relativa a:

\begin{lstlisting}
governed project baseline
+ pinned source revision
+ BA contract revision set
+ declared shared scope
\end{lstlisting}

Envelope minimo:

\begin{lstlisting}
baBaselineKey
governedBaselineKey
sourceRepositoryRevision
baContractRevisionSet
acceptedBAArtifactRef
declaredSharedScope
regressionEvidenceRef
unresolvedDiagnosticRef
reopenRule
\end{lstlisting}

Gates:

\begin{enumerate}
\item authority;
\item identity/proposition coherence;
\item provenance/source drill-down;
\item explicit uncertainty;
\item declared-scope semantic coverage;
\item post-BA regression;
\item projection readiness;
\item no unresolved contract pressure;
\item immutable baseline/change path.
\end{enumerate}

\section{Accepted Facial Access BA}

La baseline analitica corrente è:

\begin{lstlisting}
FACIAL-ACCESS-BA-R24-R1
\end{lstlisting}

derivata da:

\begin{lstlisting}
FACIAL-ACCESS-GOV-R2
source revision:
8af2257a1df94fa5a83d4853ed0a1eb4d020c429
\end{lstlisting}

\subsection{Referent principali}

\begin{lstlisting}
ControlledAreaAccess
AccessAttempt
AccessAuthorizationManagement
IdentityDetermination
PersonPresentAtAccessPoint
IdentityDeterminationRequest
RecognitionCapture
RecognitionCaptureAcquisition
RecognitionCaptureDelivery
CameraSubsystem
RecognitionProcessor
IdentityDeterminationOutcome
GovernedIdentity
AccessAuthorizationState
RequiredAccessAuthorizationCondition
AccessDecision
ConnectivityService
Confidentiality
Integrity
AuthorizedProvenance
\end{lstlisting}

\subsection{Flow minimo}

\begin{lstlisting}
PersonPresentAtAccessPoint
    -> IdentityDeterminationRequest

CameraSubsystem
    -> RecognitionCapture
       bound to IdentityDeterminationRequest

RecognitionCaptureDelivery
    CameraSubsystem -> RecognitionProcessor

RecognitionProcessor
    -> IdentityDeterminationOutcome
       SUCCESS | NEGATIVE | INCONCLUSIVE

SUCCESS
    -> GovernedIdentity X

AccessAuthorizationManagement
    -> AccessAuthorizationState concerning X

ControlledAreaAccess
    -> AccessDecision for AccessAttempt
\end{lstlisting}

\subsection{Security scope}

\begin{lstlisting}
RecognitionCaptureDelivery
    -> Confidentiality
    -> Integrity
    -> AuthorizedProvenance
\end{lstlisting}

Non:

\begin{lstlisting}
whole system -> secure
AccessDecision -> Confidentiality
Outcome -> AuthorizedProvenance
\end{lstlisting}

\section{Worked proposition set}

Un sottoinsieme rappresentativo:

\begin{lstlisting}
P-07
correlate
  correlatedItem     -> RecognitionCapture
  correlationContext -> IdentityDeterminationRequest

P-08
transfer
  behavior    -> RecognitionCaptureDelivery
  source      -> CameraSubsystem
  destination -> RecognitionProcessor
  content     -> RecognitionCapture

P-09
consumeService
  consumer -> IdentityDeterminationRealization
  service  -> ConnectivityService

P-12
constrain
  constraintTarget -> RecognitionCaptureDelivery
  constraintValue
    property   -> interconnectionMedium
    vocabulary -> [WIRED_ETHERNET]

P-13/P-14/P-15
constrain
  constraintTarget -> RecognitionCaptureDelivery
  constraintValue  -> C / I / AuthorizedProvenance

P-17
reference
  referencer -> AccessAuthorizationState
  referenced -> GovernedIdentity

P-23
decisionRule
  actor  -> ControlledAreaAccess
  input  -> IdentityDeterminationOutcome
  input  -> AccessAuthorizationState
  result -> AccessDecision
\end{lstlisting}

La lista completa e la source/provenance matrix sono nell'artifact \code{R24_FACIAL_ACCESS_BASE_ANALYSIS_R1.md}.

\section{Controlled usefulness test}

La BA governa:

\begin{lstlisting}
RecognitionCaptureDelivery.interconnectionMedium
    = WIRED_ETHERNET
\end{lstlisting}

Un consumer può escludere una minaccia la cui precondizione necessaria sia:

\begin{lstlisting}
this governed delivery uses RF/wireless medium
\end{lstlisting}

Non può inferire:

\begin{lstlisting}
the consumed provider infrastructure
is wired end-to-end
\end{lstlisting}

Questa differenza mostra perché la BA deve essere abbastanza precisa da preservare una distinzione materiale, ma non deve inventare topologia.

\section{Regression procedure}

Dopo una material BA change:

\begin{enumerate}
\item ricostruire/aggiornare la BA contro la source pin;
\item verificare ogni proposition materiale;
\item verificare diagnostics e forbidden inferences;
\item controllare property scope;
\item controllare cross-MR identity/request binding;
\item rieseguire usefulness/projection check;
\item cercare un nuovo counterexample;
\item riaprire solo il contratto minimo se necessario.
\end{enumerate}

\section{Change e revalidation}

Quando la source cambia:

\begin{lstlisting}
source change
    -> revalidationContext
    -> impacted BA elements
    -> RETAIN / REPLACE / RETIRE
    -> new accepted BA baseline
    -> rebuild projections
\end{lstlisting}

Una baseline BA accettata non viene silenziosamente mutata e riutilizzata come se le analisi storiche avessero usato la nuova semantica.

\section{Cosa non è la Base Analysis}

La BA non è:

\begin{itemize}
\item un DFD obbligatorio;
\item un modello di deployment completo;
\item una tassonomia STRIDE;
\item una sostituzione di SysML/ArchiMate/AADL;
\item una inference engine che completa il progetto;
\item una raccolta di security mechanisms;
\item la project authority;
\item un database/grafo obbligatorio.
\end{itemize}

\section{Checklist operativa}

Prima di accettare una BA baseline:

\begin{enumerate}
\item baseline authority e source revision sono pinned?
\item ogni referent riusabile ha identity stabile?
\item ogni proposition usa operator/role canonici?
\item sto usando una proposition come proxy di un comportamento?
\item ho introdotto un fatto che la source non governa?
\item un \code{NOT SPECIFIED} è stato trasformato in default?
\item ogni elemento ha provenance sufficiente?
\item le proprietà locali restano scoped al target corretto?
\item cross-request / cross-identity bindings sono preservati?
\item la regression completa passa?
\item una projection controllata consuma la BA senza reparsing del prose?
\item esiste un counterexample corrente che richiede reopen?
\end{enumerate}

\section{Stato corrente}

Per lo scope R24 esercitato dal corpus Facial Access:

\begin{lstlisting}
BA0 PASS
BA1 PASS
BA2 R3 PASS
BA3 PASS
BA4 PASS
BA5 PASS
BA6 PASS

accepted BA baseline:
    FACIAL-ACCESS-BA-R24-R1

universal completeness:
    NOT CLAIMED

reopen:
    concrete governed counterexample
    or governed source change
\end{lstlisting}

Questa guida R2 è la guida operativa human-readable corrente per la Base Analysis.

Il passo successivo è usare la documentazione e la BA stabilizzate per revisionare i capitoli di tesi completabili prima della valutazione STRIDE/STRIDE-AI.
\end{document}
